\documentclass[12pt]{article}
\usepackage[T1]{fontenc}
\usepackage[utf8]{inputenc}
\usepackage{lmodern}
\usepackage{textcomp}
\usepackage{enumitem}
\usepackage{geometry}
\geometry{margin=1in}
\begin{document}
\begin{center}
Answer Key - Psychology - Undergraduate Level Assessment 25 \
\small (Answers are ordered exactly as in the question paper.)
\end{center}

\section*{Multiple Choice}
\begin{enumerate}[label=\arabic*.]
\item \textbf{Answer:} A
\\ \textit{Options:} A) sour grapes rationalization; B) sweet lemons rationalization; C) sublimation; D) rationalization; E) projection
\\ \textit{Note:} The concept of sour grapes rationalization applies here as it describes the psychological defense mechanism where individuals attribute their failures or misfortunes to others, in this case, Jews and minorities, whom the Nazis unjustly blamed for Germany's economic and social issues during World War II.
\item \textbf{Answer:} E
\\ \textit{Options:} A) agreeableness; B) emotional stability; C) assertiveness; D) neuroticism; E) openness
\\ \textit{Note:} Individuals who score high on openness tend to be more imaginative, curious, and willing to explore new experiences, which are characteristics that facilitate the susceptibility to hypnosis.
\item \textbf{Answer:} B
\\ \textit{Options:} A) duration of the sound; B) amplitude of the wave; C) source of the sound; D) transduction of the tone; E) pitch of the tone
\\ \textit{Note:} Turning up the volume on a music player increases the amplitude of the sound wave, which corresponds to a higher intensity of sound perceived by the listener.
\item \textbf{Answer:} A
\\ \textit{Options:} A) aggressive acts can only be predicted by physical characteristics; B) feelings of anonymity lead to more uncharacteristic violence; C) aggressive acts increase the likelihood of further aggression; D) frustration is a prerequisite for any aggressive act; E) catharsis theory states that high stress levels lead to more aggression
\\ \textit{Note:} The correct answer reflects the idea that the Catharsis Theory, which suggests that expressing aggression can lead to the release of pent-up emotions, is not supported by research, indicating that aggression is more closely linked to innate physical characteristics rather than being solely a response to frustration.
\item \textbf{Answer:} A
\\ \textit{Options:} A) Minnesota Multiphasic Personality Inventory; B) Rey 15-Item Memory Test; C) Beck Depression Inventory; D) The General Aptitude Test Battery; E) Recognition Memory Test
\\ \textit{Note:} The Minnesota Multiphasic Personality Inventory (MMPI) is primarily designed to assess personality structure and psychopathology rather than specifically measuring symptom validity or detecting malingering.
\end{enumerate}

\section*{True / False}
\begin{enumerate}[label=\arabic*.]
\setcounter{enumi}{5}
\item \textbf{Answer:} False
\\ \textit{Options:} A) True; B) False
\\ \textit{Note:} The medulla oblongata primarily controls autonomic functions such as breathing and heart rate, and while it plays a role in overall brain function, it is not directly responsible for vision, which is primarily managed by the occipital lobe and other related structures.
\item \textbf{Answer:} True
\\ \textit{Options:} A) True; B) False
\\ \textit{Note:} The statement is true because the iris is a muscle that adjusts the size of the pupil, thereby regulating the amount of light that enters the eye and ultimately reaches the retina for visual processing.
\item \textbf{Answer:} False
\\ \textit{Options:} A) True; B) False
\\ \textit{Note:} A performance review serves not only to assess employee performance but also to clarify job requirements, set expectations, and facilitate employee development, making the statement False.
\item \textbf{Answer:} False
\\ \textit{Options:} A) True; B) False
\\ \textit{Note:} Crystallized intelligence encompasses knowledge and skills acquired through experience and education, which are heavily influenced by cultural factors, contrary to the idea that it is solely based on innate abilities.
\item \textbf{Answer:} False
\\ \textit{Options:} A) True; B) False
\\ \textit{Note:} EEGs with alpha and beta waves are primarily associated with relaxed wakefulness and light sleep, while stage 3 sleep is characterized by delta waves, indicating deep sleep.
\end{enumerate}

\section*{Long Form Response}
\begin{enumerate}[label=\arabic*.]
\setcounter{enumi}{10}
\item \textbf{Answer:} Convergent thinking and divergent thinking are two distinct cognitive processes with significant implications for problem-solving and creativity in psychological contexts. Convergent thinking focuses on identifying a single, correct solution to a problem, often emphasizing analytical reasoning and logic. In contrast, divergent thinking involves the generation of multiple possible solutions, encouraging creativity and innovative approaches. The effectiveness of convergent thinking is commonly seen in situations requiring precise answers, such as standardized testing or technical problem-solving, whereas divergent thinking is crucial in creative fields, fostering brainstorming and idea generation. Understanding these differences can enhance educational strategies and therapeutic approaches by promoting a balance between structured problem-solving and creative exploration.
\\ \textit{Note:} The answer correctly distinguishes between convergent thinking, which seeks a single correct solution through analytical logic, and divergent thinking, which fosters creativity by generating multiple potential solutions, highlighting their respective roles in problem-solving and creativity within psychological contexts.
\item \textbf{Answer:} Paranoid disorder is characterized by persistent, persecutory delusions, where individuals believe they are being targeted or harmed by others, or delusions of jealousy, where they suspect infidelity without evidence. There are three distinct types of paranoid disorders: paranoia, which involves a pervasive mistrust and suspicion of others; shared paranoia, where two or more individuals share the same delusional beliefs; and paranoid state, which is a temporary period of intense paranoia that may be triggered by stress or substance use. Each type presents unique challenges for individuals, affecting their social relationships, occupational functioning, and overall quality of life. Understanding these distinctions is crucial for effective treatment and support for those affected by paranoid disorders.
\\ \textit{Note:} The answer correctly identifies the primary symptoms of paranoid disorder, such as persecutory delusions and jealousy, and distinguishes between the three types of the condition-paranoia, shared paranoia, and paranoid state-while outlining their characteristics and implications for those affected.
\end{enumerate}

\vfill
\noindent\textit{End of Answer Key}
\end{document}